\documentclass[a4paper,11pt]{article}

% --- Idioma y codificación ---
\usepackage[utf8]{inputenc}
\usepackage[T1]{fontenc}
\usepackage[spanish,es-nodecimaldot]{babel}

% --- Matemática ---
\usepackage{amsmath}
\usepackage{amsfonts}
\usepackage{amssymb}
\usepackage{breqn}

% --- Listas y tablas ---
\usepackage{enumitem}
\usepackage{multirow}
\usepackage{array}

% --- Gráficos y colores ---
\usepackage{graphicx}
\usepackage{xcolor}
\usepackage{pgfplots}
\usepackage{tikz}

\usetikzlibrary{datavisualization.formats.functions}
\usetikzlibrary{calc}
\usepgfplotslibrary{polar}
\usepgflibrary{shapes.geometric}

\pgfplotsset{compat=1.18}

% --- Márgenes ---
\usepackage[
    left=2.50cm,
    right=2.50cm,
    top=2.30cm,
    bottom=2.20cm,
    headheight=15pt,
    headsep=0.45cm,
    footskip=0.90cm
]{geometry}

% --- Otros paquetes ---
\usepackage{natbib}
\usepackage{float}
\usepackage{chngcntr}
\usepackage{blindtext}
\usepackage{verbatim}
\usepackage{ifoddpage}
\usepackage{needspace}

% --- Colores ---
\colorlet{colormarco}{teal}
\colorlet{fondoseccion}{teal!12}
% --- Hipervínculos ---
\usepackage{hyperref}

\hypersetup{
    colorlinks=true,
    linkcolor=teal,
    urlcolor=teal,
    citecolor=green,
    pdftitle={Trabajo práctico 1: Matemática Aplicada 1},
    pdfauthor={Juan Andrés Cabral}
}

% --- Encabezados y pies de página ---
\usepackage{fancyhdr}

\newcommand{\encabezadoguia}{%
    \fancyhf{}
    \fancyhead[L]{%
        \small
        \href{https://cabraljuan.github.io}{Juan Andrés Cabral}%
    }
    \fancyhead[R]{\small\itshape Trabajo práctico 1}
    \fancyfoot[R]{\small\thepage}
    \renewcommand{\headrulewidth}{0pt}
    \renewcommand{\footrulewidth}{0pt}
}

\pagestyle{fancy}
\encabezadoguia

\fancypagestyle{plain}{%
    \encabezadoguia
}

% --- Marco alrededor de todas las páginas ---
% Modificar 0.75cm para acercar o alejar el marco del borde.
\AddToHook{shipout/background}{%
    \begin{tikzpicture}[remember picture,overlay]
        \draw[
            draw=colormarco,
            line width=0.6pt
        ]
        ([xshift=0.75cm,yshift=-0.75cm]current page.north west)
        rectangle
        ([xshift=-0.75cm,yshift=0.75cm]current page.south east);
    \end{tikzpicture}%
}

% --- Párrafos ---
\setlength{\parindent}{24pt}

% --- Títulos de los apartados ---
\newcommand{\seccionguia}[1]{%
    \par\addvspace{0.45cm}
    \Needspace{4\baselineskip}
    \noindent
    \colorbox{fondoseccion}{%
        \parbox{\dimexpr\linewidth-2\fboxsep\relax}{%
            \raggedright
            \strut\textbf{#1}\strut
        }%
    }%
    \par\nobreak\vspace{0.15cm}
}

% --- Ejercicios ---
\newlist{ejercicios}{enumerate}{1}

\setlist[ejercicios]{
    label=\arabic*.,
    ref=\arabic*,
    leftmargin=*,
    labelsep=0.5em,
    itemsep=0.45cm,
    topsep=0.15cm,
    parsep=0pt,
    partopsep=0pt
}

% --- Incisos ---
\newlist{incisos}{enumerate}{1}

\setlist[incisos]{
    label=\alph*),
    ref=\alph*,
    leftmargin=*,
    labelsep=0.5em,
    itemsep=0.10cm,
    topsep=0.12cm,
    parsep=0pt,
    partopsep=0pt
}

% --- Respuestas ---
% Numeración independiente de la lista de ejercicios.
\newlist{respuestas}{enumerate}{1}

\setlist[respuestas]{
    label=\arabic*.,
    ref=\arabic*,
    leftmargin=*,
    labelsep=0.5em,
    itemsep=0.35cm,
    topsep=0.15cm,
    parsep=0pt,
    partopsep=0pt
}

% --- Datos del documento ---
\title{Trabajo práctico 1: Matemática Aplicada 1}
\author{Cátedra: García Venturini - Profesor: Juan Andrés Cabral}
\date{}

\begin{document}

% ==================================================
% Título de la primera página
% ==================================================

{\centering
    {\Large\bfseries Matemática Aplicada 1\par}
    \vspace{0.12cm}
    {\normalsize Cátedra: García Venturini\par}
    \vspace{0.05cm}
    {\normalsize Profesor: Juan Andrés Cabral\par}
    \vspace{0.20cm}
    {\large\bfseries Trabajo práctico 1\par}
}

\vspace{0.20cm}
{\color{colormarco}\hrule height 0.6pt}

% ==================================================
% 1. Funciones de varias variables y curvas de nivel
% ==================================================

\seccionguia{Funciones de varias variables y curvas de nivel}

% Quitar los signos % del bloque al agregar los ejercicios.
% La opción series=guia inicia la numeración continua.
%
% \begin{ejercicios}[series=guia]
%
%     \item Escribir aquí la consigna del primer ejercicio
%
%     \begin{incisos}
%         \item Escribir aquí el primer inciso
%         \item Escribir aquí el segundo inciso
%     \end{incisos}
%
%     \item Escribir aquí la consigna del segundo ejercicio
%
% \end{ejercicios}

% ==================================================
% 2. Aplicaciones económicas
% ==================================================

\seccionguia{Aplicaciones económicas de las funciones y sus curvas de nivel}
\begin{ejercicios}[resume=guia]

    \item Considere la siguiente función de utilidad, donde \(x\) e \(y\)
    representan las cantidades consumidas de dos bienes:
    \[
    u(x,y)=x^{1/2}y^{1/2},
    \qquad x>0,\quad y>0
    \]

    \begin{incisos}
        \item Graficar las curvas de indiferencia para los niveles
        \(u=1\), \(u=2\) y \(u=3\).

        \item ¿Las curvas de indiferencia cumplen las condiciones
        deseables discutidas en clase? Justificar la respuesta.
    \end{incisos}

\end{ejercicios}

% La opción resume=guia continúa la numeración anterior.
%
% \begin{ejercicios}[resume=guia]
%
%     \item Escribir aquí la consigna del siguiente ejercicio
%
%     \begin{incisos}
%         \item Escribir aquí el primer inciso
%         \item Escribir aquí el segundo inciso
%     \end{incisos}
%
% \end{ejercicios}

% ==================================================
% 3. Derivadas parciales y diferencial
% ==================================================

\seccionguia{Derivadas parciales y diferencial}

% \begin{ejercicios}[resume=guia]
%
%     \item Escribir aquí la consigna del siguiente ejercicio
%
%     \begin{incisos}
%         \item Escribir aquí el primer inciso
%         \item Escribir aquí el segundo inciso
%     \end{incisos}
%
% \end{ejercicios}

% ==================================================
% 4. Aplicaciones económicas de las derivadas
% ==================================================

\seccionguia{Aplicaciones económicas de las derivadas y diferencial}

% \begin{ejercicios}[resume=guia]
%
%     \item Escribir aquí la consigna del siguiente ejercicio
%
%     \begin{incisos}
%         \item Escribir aquí el primer inciso
%         \item Escribir aquí el segundo inciso
%     \end{incisos}
%
% \end{ejercicios}

% ==================================================
% Respuestas
% ==================================================

\newpage
\section*{Respuestas}
\addcontentsline{toc}{section}{Respuestas}

% Las respuestas comienzan nuevamente desde el número 1.
% Usar un \item por ejercicio, respetando el orden de la guía.
%
% \begin{respuestas}
%
%     \item
%     \begin{incisos}
%         \item {\color{teal}Respuesta al primer inciso}
%         \item {\color{teal}Respuesta al segundo inciso}
%     \end{incisos}
%
%     \item
%     {\color{teal}
%         Respuesta al segundo ejercicio
%     }
%
% \end{respuestas}

\end{document}